\documentclass[12pt,letterpaper]{article}
\usepackage[utf8]{inputenc}
\usepackage[spanish]{babel}
\usepackage{amsmath}
\usepackage{graphicx}
\usepackage{booktabs}
\usepackage{caption}
\usepackage[left=3cm,right=2.5cm,top=2.5cm,bottom=2.5cm]{geometry}

\begin{document}

\begin{titlepage}
\centering
{\Large Universidad Nacional\par}
\vspace{1cm}
{\large Facultad de Ciencias --- Departamento de Química\par}
\vspace{3cm}
{\huge\bfseries Informe de Laboratorio N.\textordmasculine\ 2\par}
\vspace{0.5cm}
{\Large Titulación ácido--base\par}
\vspace{3cm}
{\large Camila Restrepo \quad Andrés Gómez\par}
\vfill
{\large Bogotá, 2026\par}
\end{titlepage}

\section{Objetivos}
\begin{itemize}
  \item Determinar la concentración de una solución de NaOH mediante titulación con
        ftalato ácido de potasio (KHP).
  \item Aplicar el concepto de punto de equivalencia.
\end{itemize}

\section{Marco teórico}
La reacción de neutralización es
\[
\mathrm{KHC_8H_4O_4} + \mathrm{NaOH} \longrightarrow \mathrm{KNaC_8H_4O_4} + \mathrm{H_2O}.
\]
En el punto de equivalencia, $n_{\mathrm{NaOH}} = n_{\mathrm{KHP}}$, por lo tanto
\begin{equation}
C_{\mathrm{NaOH}} = \frac{m_{\mathrm{KHP}}}{M_{\mathrm{KHP}}\, V_{\mathrm{NaOH}}}.
\end{equation}

\section{Resultados}
\begin{table}[h]
\centering
\caption{Datos de la titulación.}
\begin{tabular}{cccc}
\toprule
Ensayo & $m_{\mathrm{KHP}}$ (g) & $V_{\mathrm{NaOH}}$ (mL) & $C_{\mathrm{NaOH}}$ (M) \\
\midrule
1 & 0,5102 & 24,8 & 0,1007 \\
2 & 0,5087 & 24,6 & 0,1012 \\
3 & 0,5121 & 24,9 & 0,1007 \\
\midrule
\multicolumn{3}{r}{Promedio} & 0,1009 \\
\bottomrule
\end{tabular}
\end{table}

La desviación estándar de las tres determinaciones fue de $2{,}9\times10^{-4}$~M.

\section{Conclusiones}
La concentración de la solución de NaOH es $(0{,}1009 \pm 0{,}0003)$~M. El error relativo
respecto al valor nominal de 0,1~M es del 0,9\,\%.

\end{document}
